% !TeX root = ../papers/nc-p3.tex
\section{三次錐の母線上の核と標識された頂点の排除}
\label{sec:marked-cubic-cone-exclusion}
\begin{theorem}\label{thm:marked-cubic-cone-exclusion}
定理\ref{thm:low-degree-octic-spinor}の五次と三次への分解で、三次因子が印で三重になることはない。
従って三次因子は正規で、定理\ref{thm:cubic-octic-tangent-planes}の三配置に限られる。
\end{theorem}
\begin{proof}[母線が選ぶ接三次の因子]
三重の印 $p$ を頂点とする三次錐を仮定し、その底を $C$ とする。これは例外平面の接三次曲線で、
滑らかまたは既約節点三次である。出力を合同変換して $N(p)=y_0y_0^t$ と置く。
一般母線上では頂点だけが印なので、核写像の次数は $3-1=2$ である。
実際吹上げで $2L=2(3H-E)$ を真変換した母線に制限する。その原始二次核ベクトルは三空間 $W$ を張る。
直線を像とすればその固定二空間 $U$ に対する $N|_{U\to k^{4*}}$ の二列は、
原始二次核の二成分と直交する二次ベクトルにより固定像直線 $kv$ を持つ。
対称性から $v|_U=0$（比例関係を与える一次汎関数は母線上を動く）。
一方 $p$ では二つの原始二次成分が同時に消えず二列も全零にならないため
$v$ は $y_0$ に比例し、階数一対称形式の制限から $v|_U\ne0$。これは矛盾する。

$N|_W$ は母線の全域で階数1以下である。一般階数2なら三行の余因子は次数4で、
原始二次核の平方の非零定数倍であり、階数1以下の $p$ で消え得ない。
この $W$ は $W_0=\ker y_0$ とは異なる。$W_0$ への制限の $p$ での一次部は
方向 $z\in C$ について階数2の三行 $A(z)$ だからである。
基底の正規化上で $W$ は射影的に全点に延長され、そこでも制限階数条件を保ち、
同じ理由で $W\ne W_0$ である。$W\cap W_0$ 上では $A(z)$ が恒等零である。
実際 $W$ 上の階数一条件の一次部を $p$ の単元成分で取ればよい。
従って底の正規化上で $A(z)$ の二つの一次因子を順序付けられる。
両者は至るところ異なり、係数写像の次数の和は3である。定数因子は不可能で、
全ての $A(z)$ が一つの一次式の倍になると、$C$ が入力平面を張るため全平面で行列式が零になる。
滑らかな種数一曲線は次数1の非定数射影写像を持たず、滑らかな底が除かれる。

従って正規化は $\mathbb P^1$、二因子の次数は $l,3-l$（$l=1,2$）。
$W$ の法線の $W_0^*$ への射影が第一の因子なので、法線は基点のない次数 $l$ の四列 $a(z)$ である。
母線の有限座標を $w$ とすると、入力の三座標は基点なしの二項三次式、頂点の座標が $w$ であり、
$w$ は底の $\mathcal O(3)$ の値として変換する。以下、対称テンソルを出力上の二次形式の積で表す。
正規化された錐上で
\[
                 N=(w y_0+b(z))^2+a(z)(w c(z)+d(z))                 \tag{*}
\]
となる。ここで $b,c,d$ の次数は $3,3-l,6-l$、$a,c$ の $W_0$ 上の制限は各 $z$ で独立である。
実際 $W$ 上の $w$ の一次係数は $2y_0 b$ と一意に書ける（$W\cap W_0$ 上で零）。
この $b$ は $H^1(\mathbb P^1,\mathcal O(3-l))=0$ により四空間上へ持ち上がる。
一般母線で $b|_W$ は $y_0|_W$ に比例しない。比例すれば、階数一の二次式行列の定数項も
$y_0^2$ に比例するか、あるいは一次項まで比例し定数項だけが別方向になるが、後者は階数一条件の
$w^2$ 係数に反する。詳しくは $y_0|_W=(1,0,0)$ で定数項の下二行二列が零、
さらに定数項の階数が1以下なので交差成分も零となる。この場合一般核は固定二空間 $W\cap W_0$ に入り矛盾する。
比例しないとき階数一の多項式行列の原始表示から $W$ 上の形式は $(w y_0+b)^2$ に一致する
（原始一次列の平方で、先頭と一次係数を比較する）。従ってその差は $W$ の方程式 $a$ で割れ、
次数比較が $(*)$ を与える。$W_0$ での一次係数 $a c$ が相異なる二因子であることが最後の独立性である。
\end{proof}
\begin{proof}[節点錐の頂点における平面被覆]
$n=|S\cap X|$、$b_0$ を頂点以外の二重の印の個数、$u=n+b_0$ と置く。
$18\le m=u+2\le22$。次に $u=17,18$ に絞る。
正規化した錐を解消した $\mathbb F_3$ では核線束の二捩れの差が消える。
正確には他の印の逆像を吹き上げた共通解消で $L=P$ である。従って核座標は
$H=C_0+3f$ に対する $3H$ の切断として頂点切断 $C_0$ に沿い一回消える。
頂点イデアルの引戻しは $\mathcal O(-C_0)$ であり、そこの一次係数が与える核曲線は
元の節点三次の閉埋込みを正規化したもので、非共線で次数3である。

底の三つの二項三次座標に第四座標 $w^3$ を加えた有限射 $\mathbb P^2\to S$ を用いる。
頂点以外の印が分岐平面を避けるよう第四座標を選ぶ。上の $3H$ の切断は
$\sum_{i=0}^3 w^{3(3-i)}h_{3i}(v,t)$ という九次式として持ち上がる（規格化された底上の全切断を用いる）。
頂点の逆像での基底イデアルの長さは7である。最低次の三次成分は独立な三つの基点なし二項三次を張る。
その三生成元の関係次数は4,5である。実際二変数でのHilbert--Burch分解で二関係の次数の和は9、
両方とも4以上である。従って四次全体がイデアルに入る（四次の次元は $5-3\cdot2+1=0$）ので、
次数6以上の高次項は局所長 $1+2+3+1=7$ を変えない。中山の補題で同じことが従う。
他の印の逆像は各滑らかな枝で単純で、節点母線の二重の印は二枝を持つ。
よって基底長は $s=3(n-1+b_0)+7=3u+4$。平面上では $H=3h$ なので前節の直交三束の議論で
$c_2(E)=54-s=50-3u\equiv0,3\pmod4$。これと $16\le u\le20$ が所望の二値を与える。
\end{proof}
\begin{proof}[母線上の平方根の残余を分離する]
$k=3-l$ とし $t$ を $\mathcal O_{\mathbb P^1}(k)$ の値とする。$(*)$ が階数1になるには、
平方根の符号を選んで $w y_0+b+t a$ と書けることが必要十分である。
$W$ 上で二つの平方が等しく、$a$ は零ベクトルにならないからである。具体的には
\[
         w C_1=C_2,\qquad C_1=c-2t y_0,\qquad C_2=t^2a+2tb-d .
\]
$t$ をコンパクト化した $T=\mathbb F_k$ で、無限遠切断を $B$ とし、
$A=B+kf$, $J=2B+(6-l)f$ と置く。二つの列の次数は $A,J$ であり、
$C_1$ はどこでも非零、無限遠で二列は $y_0,a$ となり独立である。
$w$ を $J-A$ の値としてさらに射影直線束 $P\to T$ にコンパクト化し、
その無限遠を $B'$ とする。四式は $D=B'+J$ の四切断であり、無限遠 $B'$ には基点がない。
有限の基点は全て単純である。実際 $C_1\ne0$ で $w$ は局所的に一意、
入力と平方根を記録する写像 $(z,w,t)\mapsto(x,w y_0+b+t a)$ は各点で不分岐である。
頂点を除いた錐の各正規化枝は底の節点でもはめ込みであり、$t$ 方向の導関数 $a$ は非零である。
一方、元の入力で階数1イデアルは各印の極大イデアルである。単元成分消去後の三行の一次部が
余接空間を張ることによる。非零な平方根を記録してもこのスキームはエタールであり、
その引戻しは上の四式と一致する（非零の $a$ の積を約する）。従って全基点が被約で、
少なくとも $u-1$ 点がある。これは頂点以外の全ての正規化枝で一つ以上の符号が可能なためである。

一般の三式は一般の $[v]\in\mathbb P^3$ に対し $w C_1-C_2\in kv$ を課す。
基点以外では $v$ を通る直線 $\langle C_1,C_2\rangle$ の入射を数える。
この直線族は二次元で全三空間を覆う。実際 $z$ を固定し $t$ を動かすと、
直線は $\langle c(z),y_0\rangle$ と動点 $C_1(t)$ で交わり、一般の直線はこの固定線と異なる
（$a,c$ は $W_0$ 上で独立）。従って $t$ は直線から復元できる。
全直線族が一次元なら、$t=\infty$ の相異なる直線 $\langle y_0,a(z)\rangle$ によって尽くされて
全て $y_0$ を通ることになり、固定 $z$ の一般直線に反する。全直線は一平面に入らない。
入れば $a,c$ の $W_0$ 上の制限がともに固定二空間に入り、$A(z)$ の三行一次行列に定数核を生じる。
二次元の直線族が曲面に載るなら一般滑らかな点を通る直線が一次元以上あり、その接平面を覆うため、曲面は平面になる。
これで支配性が従う。

従って一般三式の完全交叉には基点以外の非空の被約残余 $R$ がある。
基点では一般三式は単純点の極大イデアルを生成し、他では一般のファイバーは標数零で被約である。
一般 $v$ は二つの無限遠の像を避ける。$B'$ での像は $C_1(t,z)$ の二次元以下の像であり、
$t=\infty$ の直線族の合併も二次元以下である。
射影束の交点数は
\[
          D^3=A^2+AJ+J^2=(3-l)+(6-l)+12=21-2l .
\]
ゆえに $1\le |R|\le 22-2l-u\le3$ である。これらの点は一般 $v$ に対して同じ $z$ の上で重ならない。
固定した一般 $z$ の直線族が覆う曲面は非平面的である。平面に入れば $y_0,c,a,b,d$ が全てその三空間に入り、
$(*)$ の一般母線上の核が一定となる。非平面的な一パラメータの直線族で、一般点はただ一本の母線に属する。
もし二本以上なら一般の二本が交わる（曲面の一般点から直線対への写像の像は二次元）。
互いに交わる直線族は共面または共点であり、共点族でも一般点を通るのは一本である。
固定 $z$ での例外像は高々曲線、従って全て合わせても一般 $v$ を含まない。特殊な有限個の $z$ も避けられる。

随伴により $K_P+2D$ は $T$ 上の $K_T+A+J=B+(4-l)f$ の引戻しである。
この線束の切断は残余 $R$ 上の全ての値を指定できる。実際有限 $t$ のチャートで
次数 $4-l\ge2$ の底の多項式を取れば、異なる高々三つの $z$ で補間できる。
一方、第四の切断を掛けた完全交叉の標準跡はこれらの切断を全て消し、
非空な残余上の完全跡を与える。これは全射性に反する。これで節点錐も排除した。
\end{proof}
